%%%%%%%%%%%%%%%%%%%%%%%%%%%%%%%%%%%%%%%%%%%%%%%%%%%%%%%%%%%
% --------------------------------------------------------
%  ____  _
% |  _ \| |__   ___
% | |_) | '_ \ / _ \
% |  _ <| | | | (_) |
% |_| \_\_| |_|\___/
%
% LaTeX2e Template
% Version 3.0.0 (24/02/2026)
%
% Original Author: Guillermo Jimenez (memo.notess1@gmail.com)
% License: Creative Commons CC BY 4.0
% --------------------------------------------------------
%
% =========================================================
%  AGENT TEMPLATE — READ BEFORE FILLING IN ANY SECTION
% =========================================================
%
% This file is a structured template for writing a
% scientific research paper. It is intended to be filled
% in by an AI agent reporting on a computational or
% experimental study.
%
% GENERAL WRITING GUIDELINES FOR THE AGENT:
%
%  1. AUDIENCE: Write for scientists. Assume the reader
%     has domain expertise but may not be familiar with
%     every technical detail of the implementation.
%     Explain concepts in plain, precise language.
%
%  2. NO CODE: Do not paste raw code into the paper.
%     Instead, describe algorithms, workflows, and
%     analyses in prose. The scientists reading this
%     paper have direct access to the codebase. Use
%     pseudocode or equations only when they materially
%     aid understanding of a method.
%
%  3. NARRATIVE: Find and commit to a central narrative
%     thread. What is the core question? What did you
%     discover? Every section should connect back to
%     that thread. Guide the reader from motivation
%     through findings to implications.
%
%  4. OBJECTIVITY: Report results as they are, not as
%     you might wish them to be. Do not overstate
%     performance or significance. Where results are
%     mixed, explain why.
%
%  5. TRANSPARENCY: Explicitly state limitations,
%     failure cases, assumptions, and open questions.
%     A limitation acknowledged in the paper is far
%     more credible than one left for the reader to
%     discover.
%
%  6. CITATIONS: Every factual claim about prior work
%     must be supported by a BibTeX citation. Add all
%     entries to rho.bib and cite using \cite{key}.
%     Multiple citations can be grouped: \cite{a,b,c}.
%
%  7. FIGURES AND TABLES: All figures and tables belong
%     in the Supplemental Information section (Section 5)
%     unless a result is so central that it must appear
%     inline in Results. When referencing a figure or
%     table from the main text, use \figref{} or
%     \tabref{} so that the cross-reference is
%     hyperlinked in the PDF.
%
% =========================================================
%%%%%%%%%%%%%%%%%%%%%%%%%%%%%%%%%%%%%%%%%%%%%%%%%%%%%%%%%%%

\documentclass[9pt,a4paper,twocolumn,twoside]{rho-class/rho}
\usepackage[english]{babel}

%----------------------------------------------------------
% AGENT: Fill in the fields below.
%   \doctype   — e.g., "Research Article", "Technical Report"
%   \title     — concise, informative title of the study
%   \author    — list each contributing author or system
%   \affil     — institutional affiliation for each author
%   \dates     — date the manuscript was compiled
%   \corres    — corresponding author contact information
%   \docinfo   — brief note on compilation / tools used
%   \journalname, \journal, \vol, \no — publication venue
%----------------------------------------------------------

\doctype{Research Article}
\title{[AGENT: Insert a concise, descriptive title for the study]}

%----------------------------------------------------------

\author[{\textasteriskcentered},a,1]{[Author / Agent Name]}
\author[{\textdagger},b,2]{[Supervisor / PI Name]}

%----------------------------------------------------------

\affil[a]{[Primary affiliation — department, institution]}
\affil[b]{[Supervisor affiliation — department, institution]}

%----------------------------------------------------------

\dates{This manuscript was compiled on \today}

%----------------------------------------------------------

\corres{\textsuperscript{\textasteriskcentered}Corresponding author: \href{mailto:author@institute.org}{author@institute.org}}
\corres{\textsuperscript{\textdagger}Senior / supervising author.}

%----------------------------------------------------------

\docinfo{[AGENT: Briefly note the computational environment, key software versions, or any important compilation details relevant to reproducibility.]}

%----------------------------------------------------------

\journalname{[Journal / Venue]}
\journal{[Full Journal Title, Year]}
\theday{\today}
\vol{X}
\no{X}

%----------------------------------------------------------
% ABSTRACT
%
% AGENT INSTRUCTIONS:
%   Write 150–250 words. Cover:
%     (a) the motivating problem and why it matters,
%     (b) the approach taken at a high level,
%     (c) the two or three most important findings, and
%     (d) the broader implication or next step.
%   Do not cite references in the abstract.
%   Do not use jargon without defining it.
%----------------------------------------------------------

\begin{abstract}
    [AGENT: Replace this placeholder with the abstract.
    Summarize the motivation, approach, key findings, and
    implications in 150--250 words. Be specific about
    quantitative results where possible.]
\end{abstract}

%----------------------------------------------------------

\keywords{[keyword 1], [keyword 2], [keyword 3], [keyword 4], [keyword 5]}

%----------------------------------------------------------

\begin{document}

    \maketitle
    \thispagestyle{firststyle}

%==========================================================
% SECTION 1 — INTRODUCTION
%
% TARGET LENGTH: 1–2 pages (two-column layout).
%
% AGENT INSTRUCTIONS:
%   The introduction does three things: (1) motivates the
%   problem, (2) surveys relevant prior work, and (3)
%   states what this paper contributes.
%
%   Recommended paragraph structure:
%     P1 — Hook: the scientific problem and why it is
%          important. What gap exists in current knowledge
%          or capability? Cite foundational work.
%     P2 — Prior work: what has been tried, what works,
%          and where the field falls short. Be fair and
%          thorough. Every claim needs a citation.
%     P3 — This work: what you did, at a high level.
%          Do NOT recite all results here; save that for
%          the Results section.
%     P4 — Paper outline: one sentence per section so the
%          reader knows what is coming.
%
%   Tone: engaging but precise. Find the narrative hook —
%   the reader should feel compelled to continue.
%   Avoid hype; let the science speak.
%==========================================================

\section{Introduction}

    \rhostart{[}AGENT: Begin with a compelling opening
    sentence that captures the core scientific question.
    Replace this paragraph with the motivating context
    for the study.]

    [AGENT: Prior work paragraph. Survey the relevant
    literature. Highlight what methods or findings already
    exist \cite{example_ref_1}, where they succeed, and
    where they fall short. Every factual claim must have
    a citation. Aim for 5--15 cited works.]

    [AGENT: This work paragraph. Describe the approach
    taken in this study at a high level — what was done
    and why it was the right choice given the landscape
    of prior work. Do not reproduce method details here;
    those belong in the Methods section.]

    [AGENT: Outline paragraph. Briefly state what each
    remaining section contains, e.g., ``Section~2
    presents the main results, Section~3 describes the
    methods in detail, and Section~4 provides supplemental
    analyses and extended data.'']

%==========================================================
% SECTION 2 — RESULTS
%
% TARGET LENGTH: 5–6 pages (two-column layout).
%
% AGENT INSTRUCTIONS:
%   This is the heart of the paper. Present findings
%   clearly and in a logical order — typically from
%   broad/validating results to specific/novel ones.
%
%   STRUCTURE GUIDANCE:
%     • Use one subsection per major finding or
%       analytical theme.
%     • Open each subsection with one sentence stating
%       what was found, then provide the supporting
%       evidence and context.
%     • Report numbers with appropriate precision and
%       uncertainty (e.g., mean ± std, confidence
%       intervals, p-values where applicable).
%     • Compare to baselines or benchmarks wherever
%       possible. A result only has meaning relative
%       to something.
%     • Reference figures and tables from the
%       Supplemental section using \figref{} and
%       \tabref{}. Only place a figure inline here if
%       it is the single most important result of the
%       paper and must be seen immediately.
%     • Acknowledge unexpected or negative results
%       honestly. If something did not work as
%       expected, say so and offer a hypothesis for
%       why.
%
%   NARRATIVE GUIDANCE:
%     Build a coherent story across subsections. Each
%     finding should naturally lead to the next question,
%     which the following subsection answers. By the end
%     of this section the reader should feel that the
%     central question of the paper has been addressed.
%
%   LIMITATIONS — weave these in as you report results,
%   rather than deferring all caveats to a single block
%   at the end. If a result is sensitive to a particular
%   assumption or dataset, say so in the same paragraph.
%==========================================================

\section{Results}

    %------------------------------------------------------
    % SUBSECTION 2.1 — FINDING / THEME 1
    %
    % AGENT: Replace this subsection heading with a
    % short, descriptive title for the first major
    % finding or analytical theme (e.g., "Model
    % accurately predicts X across Y conditions").
    %------------------------------------------------------
    \subsection{[First Major Finding or Analytical Theme]}

        [AGENT: State the finding in the opening sentence,
        then explain the evidence. Describe the experiment
        or analysis, the data used, the metric evaluated,
        and the result. Quantify: ``The model achieved an
        accuracy of 87.3\% (±1.2\%) on the held-out test
        set, compared to 74.1\% for the baseline.'' Cross-
        reference supporting figures or tables in the
        Supplemental section, e.g., see \figref{fig:main_result}.]

        [AGENT: If relevant, discuss what this result means
        in the context of prior work or the biological /
        scientific system under study. What does it imply?
        Note any important caveats or conditions under
        which the result holds.]

    %------------------------------------------------------
    % SUBSECTION 2.2 — FINDING / THEME 2
    %------------------------------------------------------
    \subsection{[Second Major Finding or Analytical Theme]}

        [AGENT: Follow the same structure as 2.1. State the
        finding, provide quantitative evidence, compare to a
        baseline or prior result, and note limitations. Each
        subsection should feel self-contained but connected
        to the overall narrative.]

    %------------------------------------------------------
    % SUBSECTION 2.3 — FINDING / THEME 3
    %
    % AGENT: Add or remove subsections as needed to cover
    % all major findings. Aim for 3–5 subsections total.
    % Fewer is usually better — consolidate related results
    % rather than fragmenting them into many small ones.
    %------------------------------------------------------
    \subsection{[Third Major Finding or Analytical Theme]}

        [AGENT: Continue the results narrative. By this
        subsection, the reader should be seeing the full
        picture of what was discovered. If this finding
        depends on or contrasts with an earlier one, make
        that connection explicit.]

    %------------------------------------------------------
    % SUBSECTION 2.N — SUMMARY OF RESULTS (optional)
    %
    % AGENT: Consider closing Results with a brief
    % integrative paragraph (not a subsection) that
    % synthesizes the findings, identifies the most
    % important take-away, and sets up the Discussion /
    % Conclusion if one is included.
    %------------------------------------------------------

    [AGENT (optional): Add a brief synthesis paragraph
    here that ties together the subsections above. What is
    the single most important conclusion a reader should
    draw from this Results section?]

%==========================================================
% REFERENCES
%
% AGENT INSTRUCTIONS:
%   All references are managed through BibTeX. Add every
%   entry to rho.bib following the existing format.
%   Cite using \cite{key} in the text. Multiple citations
%   may be grouped: \cite{key1,key2,key3}.
%
%   Checklist before finalizing:
%     [ ] Every factual claim about prior work is cited.
%     [ ] Every dataset or software tool used is cited.
%     [ ] All cited works appear in rho.bib with complete
%         bibliographic information (authors, year, journal
%         or venue, volume, pages or article number, DOI).
%     [ ] No reference appears in rho.bib that is not
%         cited in the text, and vice versa.
%==========================================================

\printbibliography

%==========================================================
% SECTION 3 — METHODS
%
% TARGET LENGTH: As long as needed for full reproducibility.
%   There is no page cap — err on the side of completeness.
%
% AGENT INSTRUCTIONS:
%   The Methods section must be detailed enough that an
%   expert reader could reproduce the study without access
%   to the code. Since the scientists reading this paper
%   DO have access to the codebase, you do not need to
%   transcribe code — instead, describe the logic,
%   inputs, outputs, and design choices in plain language.
%
%   Recommended subsection structure:
%     3.1 — Data / Dataset: origin, size, preprocessing
%             steps, train/val/test splits, any filters
%             or exclusions applied, and the rationale.
%     3.2 — Model / Algorithm: the architecture or
%             algorithmic approach, key hyperparameters,
%             and why these choices were made.
%     3.3 — Evaluation: metrics used, how they were
%             computed, statistical tests if any.
%     3.4 — Computational Environment: hardware, software
%             versions, and random seeds if relevant.
%
%   Writing style: use past tense ("we trained", "the
%   model was evaluated"). Be precise about quantities.
%   State assumptions explicitly. If a design choice was
%   empirically determined (e.g., via a hyperparameter
%   sweep), say so and note what was swept.
%==========================================================

\section{Methods}

    \subsection{Data}

        [AGENT: Describe the dataset(s) used. Include the
        source, how the data were collected or curated,
        the total number of samples, any preprocessing or
        filtering steps applied, and how the data were
        split for training, validation, and testing.
        Cite the original data source \cite{example_dataset}.
        State any exclusion criteria and the rationale
        behind them.]

    \subsection{Model and Approach}

        [AGENT: Describe the computational or analytical
        approach in precise but accessible language.
        Explain the key design decisions and why they were
        made. If equations help communicate the method,
        include them using the equation environment, e.g.:

        \begin{equation}\label{eq:loss}
            \mathcal{L} = [AGENT: insert loss function or
            key equation if applicable]
        \end{equation}

        Reference equations with \equref{eq:loss}. Do not
        paste code. Describe what the code does.]

    \subsection{Evaluation}

        [AGENT: Define every metric reported in the Results
        section. Explain how each metric was computed, what
        it measures, and why it is an appropriate choice
        for this study. If statistical significance tests
        were used, name the test, the significance
        threshold, and any corrections for multiple
        comparisons.]

    \subsection{Computational Environment}

        [AGENT: Specify the hardware (e.g., GPU model,
        number of GPUs, memory), operating system, and
        key software dependencies with version numbers.
        Note wall-clock training or inference times if
        they are relevant to reproducibility or practicality.
        Specify the random seed(s) used.]

%==========================================================
% SECTION 4 — SUPPLEMENTAL INFORMATION
%
% AGENT INSTRUCTIONS:
%   This section is intentionally exhaustive. Include
%   everything that supports or extends the main paper
%   but would interrupt its flow if placed inline.
%
%   What belongs here:
%     • ALL figures and tables (unless one is placed
%       inline in Results as the single central result).
%     • Extended or ablation results not in the main text.
%     • Derivations of equations used in Methods.
%     • Additional dataset statistics or visualizations.
%     • Sensitivity analyses and robustness checks.
%     • Failure case analyses.
%     • Anything the reader might want but that would
%       slow the narrative in the main text.
%
%   HOW TO ADD FIGURES:
%     \begin{figure}[H]
%       \centering
%       \includegraphics[width=0.9\columnwidth]{figures/your_figure.pdf}
%       \caption{A clear, self-contained caption that
%         explains what is shown and what the key
%         take-away is. Panels should be labeled (A),
%         (B), etc.}
%       \label{fig:your_label}
%     \end{figure}
%
%   HOW TO ADD TABLES:
%     \begin{table}[H]
%       \caption{Table title (caption goes ABOVE the table).}
%       \label{tab:your_label}
%       \begin{tabular}{lcc}
%         \toprule
%         \textbf{Column 1} & \textbf{Column 2} & \textbf{Column 3} \\
%         \midrule
%         Row value & 0.00 & 0.00 \\
%         \bottomrule
%       \end{tabular}
%       \tabletext{Note: add any clarifying notes here.}
%     \end{table}
%
%   FIGURE AND TABLE LABELING CONVENTIONS:
%     Figures: \label{fig:descriptive_name}
%     Tables:  \label{tab:descriptive_name}
%     Reference in text: \figref{fig:name}, \tabref{tab:name}
%
%   CAPTION WRITING GUIDE:
%     A good caption is self-contained — a reader should
%     understand what the figure/table shows without
%     having to read the main text. State: (1) what is
%     plotted/tabulated, (2) what the axes or columns
%     represent, (3) what the key result or pattern is.
%     Report n, error bars, and statistical tests in
%     the caption if they appear in the figure.
%
%   AGENT: Be exhaustive here. It is better to over-
%   include than to leave a reviewer wondering where
%   the supporting evidence is.
%==========================================================

\section*{Supplemental Information}
\addcontentsline{toc}{section}{Supplemental Information}

    %------------------------------------------------------
    % S1 — EXTENDED RESULTS
    %------------------------------------------------------
    \subsection*{S1 \quad Extended Results}
    \addcontentsline{toc}{subsection}{S1 Extended Results}

        [AGENT: Report any additional quantitative results
        that support the main findings but were not
        essential in the Results section. This could
        include results on additional datasets, additional
        metrics, or results for subgroups or conditions
        explored during the study.]

    %------------------------------------------------------
    % S2 — FIGURES
    %
    % AGENT: Add one \begin{figure}...\end{figure} block
    % per figure. Follow the format above. Use PDF or
    % high-resolution PNG files saved to the figures/
    % directory. Name files descriptively.
    %------------------------------------------------------
    \subsection*{S2 \quad Figures}
    \addcontentsline{toc}{subsection}{S2 Figures}

        % ------- FIGURE TEMPLATE — duplicate as needed -------
        \begin{figure}[H]
            \centering
            \includegraphics[width=0.9\columnwidth]{figures/[AGENT: filename.pdf]}
            \caption{[AGENT: Write a self-contained caption.
              (A) Describe panel A. (B) Describe panel B.
              Error bars indicate [mean ± std / 95\% CI / etc].
              $n = $ [sample size]. State the key take-away.]}
            \label{fig:fig_s1}
        \end{figure}
        % -----------------------------------------------------

    %------------------------------------------------------
    % S3 — TABLES
    %
    % AGENT: Add one \begin{table}...\end{table} block
    % per table. Caption goes ABOVE the table body.
    %------------------------------------------------------
    \subsection*{S3 \quad Tables}
    \addcontentsline{toc}{subsection}{S3 Tables}

        % ------- TABLE TEMPLATE — duplicate as needed --------
        \begin{table}[H]
            \caption{[AGENT: Table title. Describe what
              is tabulated and, if relevant, the conditions
              under which values were measured.]}
            \label{tab:tab_s1}
            \begin{tabular}{
                >{\raggedright\arraybackslash}p{0.28\columnwidth}
                >{\centering\arraybackslash}p{0.28\columnwidth}
                >{\centering\arraybackslash}p{0.28\columnwidth}
            }
            \toprule
            \textbf{[Column 1]} & \textbf{[Column 2]} & \textbf{[Column 3]} \\
            \midrule
            [Row 1 label] & [value] & [value] \\
            [Row 2 label] & [value] & [value] \\
            \bottomrule
            \end{tabular}
            \tabletext{[AGENT: Add clarifying notes, units,
              or statistical test details as needed.]}
        \end{table}
        % -----------------------------------------------------

    %------------------------------------------------------
    % S4 — ABLATIONS AND SENSITIVITY ANALYSES
    %------------------------------------------------------
    \subsection*{S4 \quad Ablations and Sensitivity Analyses}
    \addcontentsline{toc}{subsection}{S4 Ablations and Sensitivity Analyses}

        [AGENT: Describe any ablation studies (e.g., what
        happens when a key component is removed or altered)
        and any sensitivity analyses (e.g., how results
        change as a hyperparameter is varied). This is
        critical for establishing that results are robust
        and not artifacts of a specific configuration.
        Present results in tables or figures included here.]

    %------------------------------------------------------
    % S5 — FAILURE CASES AND LIMITATIONS
    %------------------------------------------------------
    \subsection*{S5 \quad Failure Cases and Limitations}
    \addcontentsline{toc}{subsection}{S5 Failure Cases and Limitations}

        [AGENT: Describe specific cases or conditions under
        which the approach fails or performs poorly. Be
        concrete — give examples. Discuss the likely
        causes and whether they are fundamental limitations
        of the approach or engineering issues that could be
        addressed with additional work. This section
        demonstrates scientific rigor and builds trust
        with the reader.]

    %------------------------------------------------------
    % S6 — ADDITIONAL METHODS DETAILS
    %------------------------------------------------------
    \subsection*{S6 \quad Additional Methods Details}
    \addcontentsline{toc}{subsection}{S6 Additional Methods Details}

        [AGENT: Elaborate on any methods that were
        summarized briefly in Section~3 but warrant more
        detail. This could include: derivations of key
        equations, descriptions of data preprocessing
        edge cases, justifications for hyperparameter
        choices, or descriptions of any manual curation
        steps. The goal is full reproducibility.]

%----------------------------------------------------------

\end{document}
